%% kapa: exam=1 round=31 year=2020 period=1 form=A subject=real_estate subject_ko=부동산학원론 seq=12 number=92
%% source: Q-Net 감정평가사 기출문제 (cst00309) artlSeq=5207219 file=extracted/1008/5207219/A형/A형/2020년도 제31회 감정평가사 1차 시험_1교시_A형.pdf
%% conversion: kordoc 4.24.51 -> tools/gen1.py (auto); review=proofread-sonnet; flags=none
\begin{question}{92}
다음과 같은 지대이론을 주장한 학자는?
\begin{framedtext}
○ 지대는 자연적 기회를 이용하는 반대급부로 토지소유자에게 지불하는 대가로 보았다. \par
○ 토지지대는 토지이용으로부터 얻는 순소득을 의미하며, 이 순소득을 잉여라고 하였다. \par
○ 토지의 몰수가 아닌 지대의 몰수라고 주장하면서 토지가치에 대한 조세 이외의 모든 조세를 철폐하자고 하였다. \par
\end{framedtext}
\begin{choices}
  \item 리카도(D. Ricardo)
  \item 알론소(W. Alonso)
  \item 헨리 조지(H. George)
  \item 마르크스(K. Marx)
  \item 튀넨(J.H. von Thünen)
\end{choices}
\end{question}
\answer{3}
